\begin{tikzpicture}[
    >=Latex,
    scale=0.85,
    every node/.style={font=\small},
]
    % ============================================================
    % PANEL A: Unsortierte Segmente
    % ============================================================
    \begin{scope}[xshift=0cm, yshift=0cm]
        \node[font=\bfseries, anchor=south] at (2.0, 4.2)
            {(a) Unsortierte Segmente};

        % 7 Segmente mit Luecken (shorten)
        \draw[gray!70, thick, ->, shorten >=3pt, shorten <=3pt]
            (0.6, 1.0) -- (1.8, 0.6);
        \draw[gray!70, thick, ->, shorten >=3pt, shorten <=3pt]
            (3.6, 2.6) -- (3.0, 3.4);
        \draw[gray!70, thick, ->, shorten >=3pt, shorten <=3pt]
            (1.8, 0.6) -- (3.2, 0.8);
        \draw[gray!70, thick, ->, shorten >=3pt, shorten <=3pt]
            (1.4, 3.2) -- (0.6, 2.4);
        \draw[gray!70, thick, ->, shorten >=3pt, shorten <=3pt]
            (3.6, 2.6) -- (3.2, 0.8);   % wird in (c) umgedreht
        \draw[gray!70, thick, ->, shorten >=3pt, shorten <=3pt]
            (3.0, 3.4) -- (1.4, 3.2);
        \draw[gray!70, thick, ->, shorten >=3pt, shorten <=3pt]
            (0.6, 2.4) -- (0.6, 1.0);

        % Endpunkt-Marker zeigen, dass es 7 isolierte Segmente sind
        \foreach \p in {(0.6, 1.0), (1.8, 0.6), (3.2, 0.8), (3.6, 2.6),
                        (3.0, 3.4), (1.4, 3.2), (0.6, 2.4)} {
            \fill[gray!90!black] \p circle (1.5pt);
        }

        \node[anchor=north, font=\scriptsize, align=center] at (2.0, -0.2)
            {7 Segmente, beliebige\\
             Speicherreihenfolge und\\
             Ausrichtung};
    \end{scope}

    % ============================================================
    % PANEL B: Wahl Startsegment + Cursor
    % ============================================================
    \begin{scope}[xshift=6cm, yshift=0cm]
        \node[font=\bfseries, anchor=south] at (2.0, 4.2)
            {(b) Startsegment};

        % Inaktive Segmente
        \draw[gray!40, thick] (3.6, 2.6) -- (3.0, 3.4);
        \draw[gray!40, thick] (1.8, 0.6) -- (3.2, 0.8);
        \draw[gray!40, thick] (1.4, 3.2) -- (0.6, 2.4);
        \draw[gray!40, thick] (3.6, 2.6) -- (3.2, 0.8);
        \draw[gray!40, thick] (3.0, 3.4) -- (1.4, 3.2);
        \draw[gray!40, thick] (0.6, 2.4) -- (0.6, 1.0);

        % Startsegment in Magenta
        \draw[magenta!80!black, very thick, ->] (0.6, 1.0) -- (1.8, 0.6);
        \fill[magenta!80!black] (0.6, 1.0) circle (3pt);
        \node[magenta!80!black, anchor=east, font=\scriptsize]
            at (0.45, 1.0) {Start};

        % Cursor am freien Ende
        \fill[red!70!black] (1.8, 0.6) circle (3.5pt);
        \draw[red!70!black, thick] (1.8, 0.6) circle (5.5pt);
        \node[red!70!black, anchor=north west, font=\scriptsize]
            at (1.95, 0.40) {Cursor};

        \node[anchor=north, font=\scriptsize, align=center] at (2.0, -0.2)
            {Beliebiges Startsegment\\
             gewählt; Cursor markiert\\
             das freie Ende};
    \end{scope}

    % ============================================================
    % PANEL C: Iteration mit Umdrehen
    % ============================================================
    \begin{scope}[xshift=0cm, yshift=-7cm]
        \node[font=\bfseries, anchor=south] at (2.0, 4.2)
            {(c) Segment umkehren};

        % Bisheriger Ringverlauf in Magenta
        \draw[magenta!80!black, very thick, ->] (0.6, 1.0) -- (1.8, 0.6);
        \draw[magenta!80!black, very thick, ->] (1.8, 0.6) -- (3.2, 0.8);

        % Cursor am aktuellen Ende
        \fill[red!70!black] (3.2, 0.8) circle (3.5pt);
        \draw[red!70!black, thick] (3.2, 0.8) circle (5.5pt);

        % Originalrichtung (gestrichelt grau, LINKS versetzt)
        \draw[gray!50, thick, dashed, ->]
            ($(3.6, 2.6)+(-0.14, 0.03)$) -- ($(3.2, 0.8)+(-0.14, 0.03)$);

        % Umgedrehte Richtung (orange, RECHTS versetzt)
        \draw[orange!80!black, very thick, ->]
            ($(3.2, 0.8)+(0.14, -0.03)$) -- ($(3.6, 2.6)+(0.14, -0.03)$);

        % Restliche Segmente grau
        \draw[gray!40, thick] (3.6, 2.6) -- (3.0, 3.4);
        \draw[gray!40, thick] (1.4, 3.2) -- (0.6, 2.4);
        \draw[gray!40, thick] (3.0, 3.4) -- (1.4, 3.2);
        \draw[gray!40, thick] (0.6, 2.4) -- (0.6, 1.0);

        % Annotation "original" - INNEN im Ring
        \node[gray!50!black, anchor=east, font=\scriptsize]
            at (2.8, 2.0) {original};
        \draw[gray!50!black, thin]
            (2.85, 2.0) to[bend right=10] (3.25, 1.7);

        \node[orange!80!black, anchor=west, font=\scriptsize]
            at (3.9, 2.4) {umkehren};
        \draw[orange!80!black, ->, thick]
            (4.0, 1.5) -- (4.2, 1.5) arc[start angle=-90, end angle=90, radius=0.25] -- (4.0, 2.0);

        \node[anchor=north, font=\scriptsize, align=center] at (2.0, -0.2)
            {Näherer Endpunkt liegt am\\
             hinteren Ende $\Rightarrow$ Segment\\
             vor Anfügen umkehren};
    \end{scope}

    % ============================================================
    % PANEL D: Geschlossener Ring + Snap-Tolerance
    % ============================================================
    \begin{scope}[xshift=6cm, yshift=-7cm]
        \node[font=\bfseries, anchor=south] at (2.0, 4.2)
            {(d) Geschlossener Ring};

        % Startpunkt-Marker
        \fill[magenta!80!black] (0.6, 1.0) circle (3pt);

        % Toleranzkreis um Startpunkt
        \draw[red!70!black, thick, dashed] (0.6, 1.0) circle (11pt);

        \node[red!70!black, anchor=east, font=\scriptsize, align=center]
            at (0.2, 0.2) {Snap-Tolerance\\$\varepsilon$};
        \draw[red!70!black, ->, thin]
            (-0.25, 0.7) to[out=45, in=200] (0.25, 0.9);

        \node[anchor=north, font=\scriptsize, align=center] at (2.0, -0.2)
            {Cursor in $\varepsilon$-Umgebung\\
             des Startpunkts\\
             $\Rightarrow$ Ring schließt};

        % Vollstaendiger Ring in Magenta
        \draw[magenta!80!black, very thick, ->] (0.6, 1.0) -- (1.8, 0.6);
        \draw[magenta!80!black, very thick, ->] (1.8, 0.6) -- (3.2, 0.8);
        \draw[magenta!80!black, very thick, ->] (3.2, 0.8) -- (3.6, 2.6);
        \draw[magenta!80!black, very thick, ->] (3.6, 2.6) -- (3.0, 3.4);
        \draw[magenta!80!black, very thick, ->] (3.0, 3.4) -- (1.4, 3.2);
        \draw[magenta!80!black, very thick, ->] (1.4, 3.2) -- (0.6, 2.4);
        \draw[magenta!80!black, very thick, ->] (0.6, 2.4) -- (0.5, 1.2);
    \end{scope}
\end{tikzpicture}